\vspace{-2.5mm}
\section{Related Work}
\label{sec:related}
\vspace{-1mm}

\textbf{Efficient Transformers.} 
In order to improve the inference speed and reduce the memory footprint of Transformers, multiple different approaches have been proposed.
These can be broadly categorized as follows:
(i) efficient architecture design~\cite{iandola2020squeezebert,kitaev2019reformer, lan2019albert,sun2020mobilebert, wang2020linformer,wu2020lite};
(ii) hardware-software co-design~\cite{ham20203, ham2021elsa, tambe2021edgebert, wang2021spatten};
(iii) knowledge distillation~\cite{jiao2019tinybert, sanh2019distilbert,sun2019patient, wang2020minilm};
(iv) quantization~\cite{kim2021bert, shen2020q, zadeh2020gobo, zafrir2019q8bert};
(v) neural architecture search~\cite{chen2020adabert, so2019evolved, so2021primer,  wang2020hat, xu2021bert, yin2021autotinybert};
and (vi) pruning.
In this paper, we focus on pruning and briefly discuss the related works.

\textbf{Transformers Pruning.}
Pruning has been a promising way to remove unimportant weights in neural networks.
Pruning can be largely categorized into unstructured and structured pruning. 
For unstructured pruning, magnitude-based~\cite{gale2019state}, first-order~\cite{sanh2020movement}, and second-order~\cite{kurtic2022optimal} pruning methods, and the lottery ticket hypothesis~\cite{chen2020lottery, chen2020earlybert, frankle2018lottery, prasanna2020bert} have been explored for Transformers.
While these methods can substantially compress the model size, commodity hardware such as GPUs can hardly take advantage of the unstructured sparse patterns for model inference speedup.

For this reason, a number of structured pruning methods have been introduced to remove coarse-grained sets of parameters in Transformers.
For example, to prune structured sets of parameters in weight matrices, low-rank factorization~\cite{wang2019structured}, block-wise sparsity~\cite{li2020efficient}, and tile-wise sparsity~\cite{guo2020accelerating} were studied.
Furthermore, as more coarse-grained methods, attention head pruning~\cite{michel2019sixteen, voita2019analyzing} and layer dropping~\cite{fan2019reducing, sajjad2020effect} have been popularly used.
Taking a step further, recent approaches~\cite{chen2021chasing, khetan2020schubert,  lagunas2021block, lin2020pruning, liu2021rosita, xia2022structured, yao2021mlpruning} have explored jointly pruning Transformers with different pruning granularity and principles, maximizing the model efficiency in every dimension.
Orthogonally, another thread of work~\cite{fan2019reducing, hou2020dynabert, liu2021ebert, xin2020deebert, zhou2020bert} has shown that Transformers can be dynamically pruned at inference time.

Unfortunately, while the structured pruning methods can achieve high compression rates and speedups, they are often difficult to use in practice.
One reason for this is the high computational cost of additional training during or after pruning, which can be up to 10$\times$~\cite{lagunas2021block, xia2022structured} compared to that of the original model training.
Another reason is the high complexity of the pruning pipelines~\cite{hou2020dynabert, lagunas2021block, liu2021rosita, yao2021mlpruning}, where each pruning stage often requires rewriting the training code and introduces additional hyperparameters to tune.

\textbf{Post-training Model Compression.}
Post-training compression methods have been widely studied in quantization.
These methods, categorized as post-training quantization (PTQ), perform quantization without any retraining, thereby avoiding the additional training cost and user intervention.
Multiple PTQ techniques have been proposed to effectively mitigate the accuracy degradation without retraining~\cite{banner2018post, hubara2020improving, nagel2020up, zhao2019improving}.
 
Although not as much as for quantization, post-training schemes have also been explored for unstructured~\cite{frantar2022spdy, lazarevich2021post, mussay2019data} and structured pruning of CNNs.
For structured pruning, \cite{kim2020neuron, srinivas2015data, yvinec2021red} proposed ways to group and merge similar neurons in a CNN.
However, we find it difficult to extend those techniques to pruning Transformers because they require the model to have a repeating structure of a linear layer and an element-wise nonlinearity, which is not the case for the multi-head attention layers of Transformers.
Even for the feed-forward network layers of Transformers,
\cite{kim2020neuron, srinivas2015data} can hardly be used because they rely on a certain characteristic of ReLU while many Transformers~\cite{brown2020language, devlin2018bert, yang2019xlnet} use GELU~\cite{hendrycks2016gaussian} instead of ReLU.

Motivated by the fact that the existing post-training CNN pruning techniques cannot be applied to Transformers, in this paper we propose a novel post-training pruning method with a focus on Transformers.
However, we would like to note that the underlying principles in our approach are general enough to be extended to pruning other types of model architectures such as CNNs.
